\documentclass[sn-mathphys-num]{sn-jnl}

\usepackage{graphicx}
\usepackage{multirow}
\usepackage{amsmath,amssymb,amsfonts}
\usepackage{amsthm}
\usepackage{mathrsfs}
\usepackage[title]{appendix}
\usepackage{xcolor}
\usepackage{textcomp}
\usepackage{manyfoot}
\usepackage{booktabs}
\usepackage{algorithm}
\usepackage{algorithmicx}
\usepackage{algpseudocode}
\usepackage{listings}
\usepackage{tabularx}
\usepackage{longtable}
\usepackage{adjustbox}
\graphicspath{{figures/}} 

% -------------------------------------------------------------------------
% New column type for centred tabularx columns
\newcolumntype{Y}{>{\centering\arraybackslash}X}
% -------------------------------------------------------------------------

\begin{document}

% =========================================================================
% TITLE
% =========================================================================

\title[Addressing soil pollution]{Addressing soil pollution:
    Advanced techniques for effective monitoring of soil ecosystem services}

% =========================================================================
% AUTHORS AND AFFILIATIONS
% =========================================================================

\author*[1,2,3]{\fnm{Polina} \sur{Lemenkova}}
\email{polina.lemenkova@tech-orleans.fr}
% ORCID: 0000-0002-5759-1089

\affil*[1]{%
    \orgdiv{Bureau de Recherches G\'eologiques et Mini\`eres},
    \orgname{BRGM},
    \orgaddress{%
        \street{3 Avenue Claude-Guillemin},
        \city{Orl\'eans Cedex~2},
        \postcode{45060},
        \country{France}}}

\affil[2]{%
    \orgdiv{Institut d'Etudes Avanc\'ees du Val de Loire},
    \orgname{LE STUDIUM},
    \orgaddress{%
        \street{1 Rue Dupanloup},
        \city{Orl\'eans},
        \postcode{45000},
        \country{France}}}

\affil[3]{%
    \orgname{La Technopole d'Orl\'eans},
    \orgaddress{%
        \street{1 Avenue du Champ de Mars},
        \city{Orl\'eans},
        \postcode{45100},
        \country{France}}}

% =========================================================================
% ABSTRACT AND KEYWORDS
% =========================================================================

\abstract{%
Soil pollution is a pressing global concern linked to natural hazards and
industrial activities. An integrative approach to soil monitoring is suggested
to identify primary pollutants such as metals, pesticides, microplastics, and
PFAS. This review synthesizes insights from over 100 scholarly articles and
emphasizes the necessity for innovative analytical tools and methodologies for
effective soil screening. Recent computational advancements have improved
pollutant detection, categorized into four groups: process-based, empirical,
composite, and statistical methods, with hyperspectral imagery and sensors
highlighted for their effectiveness. Nonetheless, challenges like technical
limitations and varying soil structures exist. Future research should leverage
computational modeling and machine learning for enhanced real-time monitoring
of pollutants, aiming for a sustainable management framework for soil
ecosystem services.}

\keywords{soil pollution, nature conservation, food production, nutrient
    cycling, soil contamination}

\maketitle

% =========================================================================
\section{Introduction}\label{sec1}
% =========================================================================

Soil ecosystem services (SES) focus on sustainable management of soil and
vegetation to provide essential resources for human use. Key aspects of
maintaining soil health include monitoring climate and hydrological contexts,
regulating agricultural practices, facilitating nutrient cycling and
decomposition, and identifying contamination issues. Soil contamination
primarily arises from major pollutants such as pesticides (including
insecticides, herbicides, fungicides, and bactericides), mineral oils,
microplastics, per- and polyfluoroalkyl substances (PFAS), nutrients, and
heavy metals. The increasing spread of contaminated soil poses a significant
environmental threat, necessitating the implementation of effective monitoring
strategies. Soil pollution has become a significant global environmental
issue, with recent evidence connecting it to both natural hazards and
industrial activities that increase contaminant levels. An integrative
approach to soil monitoring is proposed as a strategy for identifying primary
pollutants, which include metals, pesticides, microplastics, PFAS, and
various nutrients. This narrative integrative review synthesizes existing
literature on soil pollution monitoring techniques, drawing from a
comprehensive analysis of more than 100 scholarly articles concerning soil and
agricultural practices. These papers are sourced from reputable databases such
as Scopus, Web of Science (WoS), and Google Scholar.

The study emphasizes the need for innovative analytical tools and experimental
methodologies to enhance soil screening processes and effectively address the
challenges posed by soil contamination. Recent advancements in computational
methodologies have significantly improved the detection of soil pollutants,
crucial for enhancing ecosystem services. A systematic review has categorized
effective pollutant identification techniques into four main groups:
process-based, empirical, composite, and statistical methods. Among the
advanced strategies, hyperspectral imagery and sensors stand out for their
ability to detect contaminant patterns in soil. However, challenges remain,
such as technical limitations, varying soil structures, and the restricted
efficacy of existing instruments. Future research is recommended to
incorporate computational modeling to predict soil behavior through
high-resolution screening via remote sensing and machine learning
technologies. These innovative tools can utilize real-time data to monitor the
behavior of chemical pollutants, microplastics, and pesticides across soil and
water systems. Implementing these approaches could lead to a scalable and
sustainable framework for managing soil ecosystem services effectively.

High spatial heterogeneity in soil microstructure significantly impacts the
environmental behavior of soil elements. Soil aggregates, composed of organic
matter, minerals, and bacteria along with their interactions, play a crucial
role in the structural and functional value for agricultural crops and natural
vegetation. This highlights the critical need for consistent monitoring of
soil health, as illustrated in Figure~\ref{fig01}. The micro-scale interfaces
between these aggregates and soil pathways exhibit nonlinear characteristics,
showcasing the complexity of soil interactions. As a vital component of the
biophysical landscape, soil is influenced by various geographic and ecological
factors that shape its qualities and behaviors. These factors include the
geological and geomorphologic structure, climatic conditions, biota, and human
activities, which collectively determine the ecosystem services provided by
soil \cite{Klaucol2013,Klaucol2017,Lemenkova2023}. Therefore, in addition to
its biochemical properties, soil's significance as a landscape element is
closely tied to its natural components and their ecological interactions
\cite{Lemenkova2025a}.

During the last decade, numerous computational methods have been developed to
identify various soil components and their interactions with pollutants, which
play a crucial role in controlling the structure, processes, and behavior of
these pollutants in soils. Previous studies analyzed the application of
instrumental techniques and advanced methods, revealing significant
advancements in ecological soil monitoring. These improvements have enhanced
our understanding of how pollutants behave within organic soil aggregates.

This review systematically summarizes research focused on different types of
pollutants and the tailored methods used to detect them in soil. Emphasis is
placed on approaches that estimate how pollutants bind to minerals, organic
compounds, bacteria, and their associated complexes within soil aggregates,
all aimed at reducing pollution levels in soil horizons. The goal of this
study is two-fold: first, to classify the patterns of major soil pollutants
that affect soil structure; and second, to investigate existing methods for
revealing soil contamination caused by different pollutants while identifying
their strengths and limitations for ecological monitoring. 

This comprehensive
review focuses on modern methods for evaluating sources of soil pollution and
their impact on ecosystem services. Traditional approaches to detecting soil
pollutants are limited, prompting a review of current trends and methods in
soil ecotoxicology, chemical contamination, and ecological risk assessment.
The study also explores the statistical distribution of contaminants within
the Eurasian Economic Union (EEU) region to analyze the effects of various
contaminants, detect key sources of their inflow, and propose solutions for
soil remediation. Notably, this study represents the first attempt to present
a scoping review of soil pollution that addresses the needs of ecosystem
services, evaluates different pollutants, and summarizes the measurement
techniques and remediation methods for effective environmental monitoring.

% =========================================================================
\section{Material and Methods}\label{sec2}
% =========================================================================

Speciation and distribution of major pollutants in soil organic complexes
requires a review of their characteristics. Soil contamination is primarily
caused by hazardous substances stemming from both human and natural
activities, Figure~\ref{fig01}. 

\begin{figure}[htbp]
    \centering
    \includegraphics[width=\textwidth]{Fig_01}
    \caption{Soil structure and major sources of its pollution. Diagram source: author.}
    \label{fig02}
\end{figure}

% -------------------------------------------------------------------------
\subsection{Pollutant Groups}\label{sec2.1}
% -------------------------------------------------------------------------

Major types of soil pollution include the following groups:

\begin{itemize}
    \item \textbf{Heavy metals:} arsenic (As), cadmium (Cd), chromium (Cr),
        mercury (Hg), lead (Pb), cobalt (Co), nickel (Ni), and at higher
        concentrations---copper (Cu), zinc (Zn), iron (Fe).
    \item \textbf{Mineral oils:} petroleum-based hydrocarbons.
    \item \textbf{Organic compounds:} PAHs (Polycyclic Aromatic
        Hydrocarbons), Polychlorinated biphenyl (PCB), and Tributyltin (TBT)
        group.
    \item \textbf{Nutrients:} nitrogen (N), phosphorus (P).
    \item \textbf{POPs} (Persistent Organic Pollutants), including PFAS
        (Per- and Polyfluoroalkyl Substances).
\end{itemize}

The sources of pollutant inflow into soil ecosystems are attributed to both
anthropogenic and natural origins. Anthropogenic sources encompass
industrialization, mining activities, intensive farming utilizing agricultural
chemicals such as pesticides, irrigation with contaminated water,
urbanization, waste disposal, and incidental spills. In contrast, natural
sources arise from atmospheric deposition due to polluted air and geological
processes, such as the release of heavy metals during earthquakes. These
pollutants introduce toxic chemicals into the soil, affecting groundwater,
crops, and broader ecosystems, ultimately posing risks to human health through
contaminated food, water, and air.

% -------------------------------------------------------------------------
\subsubsection{Pollutant Group~1: Heavy Metals}\label{sec2.1.1}
% -------------------------------------------------------------------------

In Europe, the most significant and detrimental type of pollutant
affecting soil quality is heavy metals, as indicated by statistics from the
European Environment Agency (EEA), \cite{EEA2009}. Heavy metals (HM), such as arsenic (As), cadmium (Cd), chromium (Cr),
mercury (Hg), lead (Pb), cobalt (Co), and nickel (Ni), are toxic trace
elements that pose significant risks to soil and crop growth. These metals are
of particular concern due to their harmful effects on the environment and
agriculture. A schematic diagram illustrating the major sources of HM
pollution is presented in Figure~\ref{fig02}.

\begin{figure}[htbp]
    \centering
    \includegraphics[width=\textwidth]{Fig_02}
    \caption{Paths of HM pollutants into soil. Diagram source: author.}
    \label{fig02}
\end{figure}

Major sources of HM pollution include industrial and mining activities
\cite{Li2019,Wang2020,AliZerrouki2021}, with mine wastes also significantly
contributing to HM contamination \cite{Rasafi2021}. Moreover, agricultural
soil contamination occurs through the use of fertilizers, pesticides, manure,
and sewage sludge, along with the irrigation of crops with polluted water.
Urban areas face further HM pollution due to atmospheric deposition from
vehicle emissions and fossil fuel combustion \cite{Adewumi2024}. Waste
disposal practices linked to municipal and industrial garbage, along with
sewage, exacerbate HM soil contamination \cite{LindhLemenkova2022a}.

The cycle of HM transfer begins in the soil profile, where HM is transported
to the groundwater horizon and subsequently absorbed by plants and crops via
their root systems in the rhizosphere. This process ultimately leads to HM
returning to humans. Heavy metals significantly impact soil microbial
communities and enzymatic activities, serving as a primary source of
rhizosphere pollution. This pollution alters the physico-chemical properties
of soil and its microbiological dynamics. Concerns regarding chemical and HM
soil contamination arise due to their persistence and bioaccumulation within
soil ecosystems, which can be assessed through the detection of leaching.

% -------------------------------------------------------------------------
\subsubsection{Pollutant Group~2: Pesticides}\label{sec2.1.2}
% -------------------------------------------------------------------------

Pesticides are a category of pollutants characterized by their potential to
be absorbed by various organisms and their ability to drift, runoff, leach,
and evaporate into soil, water, and air. Specifically designed as chemical
compounds, pesticides serve the purpose of eliminating pests and undesired
plants. However, excessive and irresponsible use of these substances can
adversely affect soil organisms and the overall environment. This can manifest
in heightened toxicity levels, reduced biodiversity, diminished soil
fertility, and disruption to the natural processes of nutrient cycling. Key
nutrients such as nitrogen~(N), phosphorus~(P), and potassium~(K) are
particularly affected, leading to alterations in the organic matter content
within the soil \cite{Nauanova2025}. The ability of soil ecosystems to adapt
to and endure pesticide pollution is evident in their reduced resilience,
which is contingent upon the adaptation and chemical tolerance of soil
microorganisms. Furthermore, as pesticides are capable of being absorbed by
plants, they can enter the food chain, posing potential hazards to human
health \cite{GonzalezRodriguez2011}.

The cycle of pesticides within soil ecosystems is characterized by the
movement and transformation of these chemicals across various ecosystem
levels. A schematic diagram illustrating the major sources of pesticides
as pollutants of soil is presented in Figure~\ref{fig03}.

\begin{figure}[htbp]
    \centering
    \includegraphics[width=\textwidth]{Fig_03}
    \caption{Paths of pesticides into soil. Diagram source: author.}
    \label{fig03}
\end{figure}

Leaching occurs when pesticides penetrate and contaminate
groundwater, thereby elevating ecotoxicological risks associated with soil
contamination \cite{Toumi2025,PerezLucas2025}. The detrimental effects of
pesticides extend to the structural integrity of the soil, resulting in
degradation that impairs plant growth and reduces the soil's ability to absorb
and retain moisture \cite{Kelsey1968}. This degradation in turn diminishes
soil fertility, adversely affecting soil organisms including microorganisms
and earthworms \cite{Ponder2002}, and disrupts critical processes such as
microbial community dynamics and phosphorus cycling. Research indicates that
alterations in soil microbial communities further disrupt nutrient cycling
\cite{Yasir2025}. The harm inflicted by pesticides on non-target species
contributes to a decline in soil biodiversity, causing significant shifts in
the soil ecosystem's structure \cite{Marcelino2024}.

% -------------------------------------------------------------------------
\subsubsection{Pollutant Group~3: Microplastics}\label{sec2.1.3}
% -------------------------------------------------------------------------

Pervasive and persistent pollutants known as microplastics (MP) consist of
materials like polyethylene and polypropylene, and they also serve as carriers
for additional pollutants, particularly heavy metals \cite{Wang2024a}. A schematic diagram illustrating the major sources of microplastics
as pollutants of soil is presented in Figure~\ref{fig04}.

\begin{figure}[H]
    \centering
    \includegraphics[width=\textwidth]{Fig_04}
    \caption{Paths of microplastics into soil. Diagram source: author.}
    \label{fig04}
\end{figure}

The presence of microplastics affects a variety of ecosystem processes and
functions in soils, influencing physicochemical characteristics such as
porosity, enzymatic activities, microbial activities, and ultimately, plant
growth and yield \cite{Akdogan2019}. They contribute to soil degradation by
altering the physico-chemical properties, which can impact soil structure and
fertility \cite{Chen2020}. Specifically, the effects of microplastics on
microbial communities can lead to significant disruptions in nutrient cycling,
notably affecting the uptake of phosphorus~(P) and the function of acid
phosphatase \cite{Zhou2024,Wang2022}. Moreover, the absorption of
microplastics by plants is associated with stunted growth, as well as
detrimental impacts on root and seed development \cite{Carr2016}. 


The pollution caused by MP poses a significant ecological threat primarily by
affecting soil structure through alterations in water--air flow within soil
pores due to wet--dry cycles. This disturbance impairs the hydraulic
conductivity of the soil by modifying its water retention capacity, leading to
changes in density and water content in the soil aggregates. Consequently, the
composition of soil particles is also altered
\cite{Jing2023,LindhLemenkova2023d,LindhLemenkova2023e}. Additionally,
microplastics can be absorbed by plants in the rhizosphere, creating pathways
for these pollutants to enter the food chain \cite{Chen2025}. The risk
associated with MP pollution in soil ecosystems is exacerbated by a variety of
sources, including atmospheric deposition, agricultural practices that involve
sludge and plastic use, as well as industrial activities.

% -------------------------------------------------------------------------
\subsubsection{Pollutant Group~4: PFAS}\label{sec2.1.4}
% -------------------------------------------------------------------------

PFAS (per- and polyfluoroalkyl substances) are artificially created chemicals
designed to enhance product quality and are commonly found in everyday items
such as cookware, food packaging, and water-resistant clothing. Their strong
carbon--fluorine bonds contribute to their persistence in environmental
matrices, particularly in soil. A diagram showing the major sources of PFAS
into soil structure is indicated in Figure~\ref{fig05}.

\begin{figure}[H]
    \centering
    \includegraphics[width=\textwidth]{Fig_05}
    \caption{Paths of PFAS into soil. Diagram source: author.}
    \label{fig05}
\end{figure}

PFAS are frequently present alongside
microplastics (MP), creating complex interactions; for instance, PFAS can
adsorb microplastics, increasing their mobility and extending the risk of
environmental exposure. Additionally, materials and textiles that degrade into
microplastics complicate PFAS removal \cite{Shen2024}. The impact of PFAS on
soil ecosystems is multifaceted, affecting essential nutrient cycles such as
carbon~(C) and nitrogen~(N), as well as enhancing litter
decomposition---an important ecological process---and altering soil pH
levels. The ramifications of PFAS contamination extend beyond soil,
influencing both water sources and posing risks to human health through the
food and water supply
\cite{Hossini2025,Appel2022,Cheng2025,Arnesdotter2025}. For detecting soil
contamination by PFAS, various computational methods have been explored,
highlighting the ongoing need for effective monitoring strategies
\cite{Xu2023}.

% -------------------------------------------------------------------------
\subsubsection{Pollutant Group~5: Nutrients}\label{sec2.1.5}
% -------------------------------------------------------------------------

Excessive levels of major nutrients---specifically nitrogen~(N),
phosphorus~(P), and potassium~(K)---primarily stem from the application of
fertilizers and manure \cite{Baghaie2020,Luo2017,Hepperly2009}. The
accumulation of these nutrient pollutants poses significant risks to soil and
groundwater quality, ultimately leading to reduced crop yields. High
concentrations of NPK can trigger environmental issues such as algal blooms
and eutrophication, which have further implications for aquatic ecosystems
\cite{Pan2019,Guan2024,Xu2025}. A diagram showing the major sources of nutrients
into soil structure is indicated in Figure~\ref{fig06}.

\begin{figure}[H]
    \centering
    \includegraphics[width=\textwidth]{Fig_06}
    \caption{Paths of nutrients into soil structure. Diagram source: author.}
    \label{fig06}
\end{figure}

When assessing soil contamination levels caused by these nutrients, it is
crucial to develop a clear set of criteria for measuring their concentrations.
This involves a comprehensive evaluation that considers factors such as
nutrient concentration rankings in relation to soil texture types, including
sand, silt, and clay. These soil characteristics are vital, as they directly
influence porosity and permeability, which in turn affect water runoff
behavior. For instance, soils with low porosity and permeability are more
prone to runoff, exacerbating nutrient pollution in agricultural fields.
Therefore, understanding the interplay between soil properties and nutrient
concentrations is essential for managing agricultural practices and minimizing
environmental impacts. 

Nutrients can negatively impact plant development in various ways, primarily
through disrupted growth patterns and stunted root development
\cite{Gang2004,Guo2021}. Overuse of nutrients not only affects soil nutrient
balance but can also lead to alterations in soil pH, which further
complicates plant health \cite{Rahman2014}. The balance of essential minerals
absorbed by plants is crucial, and excess nutrients can hinder their
acquisition, affecting both plant and microbial growth
\cite{Ma2024,Vishwanath2025,Choudhary2022}. This imbalance has downstream
effects, as it impacts microbial activity, which is responsible for the
mineralization of organic matter and the decomposition of plant litter,
contributing to soil fertility \cite{Atoloye2024}. Additionally, nutrient
deposition from atmospheric sources, such as ammonia~(NH$_{3}$) from
cultivated areas, can harm natural terrestrial habitats, including adjacent
forests \cite{Koukianaki2025}.

% -------------------------------------------------------------------------
\subsection{Data Quantification}\label{sec2.2}
% -------------------------------------------------------------------------

Data quantification plays a crucial role in modeling soil behavior when
subjected to pollutants and chemical contaminants. This process involves the
parametrization of model dependencies, which can be quantitatively assessed
and translated into mathematical relationships based on measured values from
laboratory soil analysis. Various modeling tools are utilized in this context,
as detailed in Table~\ref{tab1}, which presents the state-of-the-art (SOTA)
techniques employed for accurate soil behavior predictions in the presence of
contaminants.

\begin{table}[htbp]
\caption{Example of SOTA techniques and instruments used for detecting
    concentration of pollutants and evaluating level of soil
    contamination.}\label{tab1}
\begin{tabularx}{\textwidth}{@{} p{1.0cm} p{2.0cm} p{5.0cm} X @{}}
\toprule
\textbf{No} & \textbf{Pollutant} & \textbf{Methods of detection}
    & \textbf{Supporting References} \\
\midrule
1 & Concentration of HM (Cd, Pb, Me)
    & Advanced methods: X-ray fluorescence (XRF) for determining the
      elemental composition of soil.
    & \cite{Lu2023,LindhLemenkova2023b,Wan2019,Mukhopadhyay2020,Adewumi2024,Feng2024,Hu2020,Li2019} \\[4pt]
2 & PAH, Tributyltin (TBT)
    & Spectroscopy and X-ray diffraction through Computed Tomography (CT)
      for non-destructive 3-D internal imaging of soil structure.
    & \cite{LindhLemenkova2022d,Deceuster2012,Racic2025,LindhLemenkova2022a,Xu2026,Stojic2019} \\[4pt]
3 & Organic pollutants (PCB)
    & Gas chromatography coupled with mass spectrometry (GC-MS) or electron
      capture detection (GC-ECD).
    & \cite{Bustamante2025,Anagnostou2025,Sandu2025,Shen2024,Appel2022,Arnesdotter2025} \\[4pt]
4 & Soil NPK levels
    & E-sensors, optical transducers, or chemical analysis for evaluation of
      NPK concentration.
    & \cite{Hossain2024,Pal2024,Ma2024,Guan2024,Atoloye2024,Luo2017} \\[4pt]
5 & pH of soil slurry
    & pH-meter or colorimetric test strip using a soil--distilled water
      slurry; logarithmic units on a scale from 0 to 14 (7 is neutral,
      $<$7 acidic, $>$7 alkaline).
    & \cite{Zireeni2023,LindhLemenkova2022b,Rahman2014,Hepperly2009,LindhLemenkova2023c,Guo2021} \\[4pt]
6 & Soil texture
    & Relative \% of sand, silt, and clay particles quantified by particle
      size analysis (soil texture triangle); e.g., 2.0--0.05\,mm diameter.
    & \cite{Polakowski2023,LindhLemenkova2023d,Mozaffari2026,LindhLemenkova2023a,LindhLemenkova2023e,Jing2023,Wang2022} \\
\bottomrule
\end{tabularx}
\end{table}

The collected dataset is utilized for modeling soil contamination, focusing on
the concentration of metals, pesticides, microplastics, PFAS, and nutrients
across various case studies in specific habitats. Analyzing data on pollutants
and biota is crucial for defining key attributes of these contaminants,
including their values and distribution. A significant challenge in data
analysis lies in automating the workflow process, which can be accomplished
using methods or scripts in R or Python
\cite{Olah2025,Lemenkova2022,Wang2024b}. This automation not only facilitates
data handling but also enhances the precision of data analysis.

Modeling plays a vital role in quantifying soil parameters and visualizing the
relationships among components within the soil ecosystem under varying
scenarios, such as climate change, the construction of mining sites, and
planned urban development. The use of Python or R libraries aids in modeling
data parameters, allowing researchers to uncover links between pollutants and
soil, as well as analyze resistivity through the creation of graphical models
and visualizations.

% =========================================================================
\section{Computational and Spectroscopic Techniques in Soil Pollution Detection}\label{sec3}
% =========================================================================

Current trends in computational methods for detecting pollutants in soil
aggregates encompass a range of traditional and advanced approaches. The
traditional methods primarily involve extraction techniques that analyze the
various forms of pollutants present in the soil \cite{Sutherland2010}. In
contrast, advanced methodologies utilize sophisticated surface analytical
techniques, such as synchrotron radiation
\cite{Hu2020,Rosskopfova2012,LindhLemenkova2023a}. Synchrotron radiation
employs advanced X-ray technologies that enable precise identification of
metal types, alongside detecting their chemical forms and speciation. These
advanced techniques are vital for understanding the reaction mechanisms of
pollutants within soil aggregates and for modeling their distribution at a
microscopic level. Instrumental technologies that exemplify this advanced
approach include Fourier transform infrared spectroscopy (FTIR) \cite{Xu2026}
and X-ray photoelectron spectroscopy (XPS)
\cite{Guo2025,LindhLemenkova2023c}, both of which significantly enhance the
capability to detect and analyze pollutants in soil systems.

The methods discussed exhibit several drawbacks primarily linked to the
detection limits of current analytical techniques, which tend to be inadequate
in the context of highly homogeneous samples. Additionally, the adsorption
percentage of contaminants on soil components is often low and uneven, posing
challenges for accurate detection. Proposed solutions to address these issues
include the development of innovative and advanced instruments for assessing
soil quality, alongside the exploration of alternative methods such as remote
sensing \cite{Fan2025,Lemenkova2025b}, machine learning
\cite{Proshad2025,Lemenkova2025d,Lemenkova2024b}, and computer vision
\cite{Feng2024}.

% -------------------------------------------------------------------------
\subsection{Computational Approaches for Identifying Pollutant Speciation
    in Soil}\label{sec3.1}
% -------------------------------------------------------------------------

The methodologies for detecting soil pollution discussed in this work focus on
the state-of-the-art (SOTA) models for soil ecosystem services (SES). The workflow summary for SES modelling related to identifying risks from soil pollution on soil ecosystem services involves a stepwise approach, which is
detailed in Table~\ref{tab3}. These
models are categorized into empirical models and process-based models, which
are essential for soil protection and remediation. Their strengths lie in
their capability to quantify and map ecosystem services, as supported by
studies from \cite{TeranGomez2025,ElSorogy2025,Lemenkova2020,Lemenkova2021,Sandu2025}.
Additionally, statistical or data-driven approaches enhance the methodologies,
as noted by \cite{Li2025a,Racic2025}.

\begin{table}[htbp]
\caption{Current modelling approaches used for the SES: strengths and
    limitations.}\label{tab3}
\begin{tabularx}{\textwidth}{@{} l p{2.5cm} p{2.5cm} p{2.5cm} X @{}}
\toprule
\textbf{Model} & \textbf{Characteristics} & \textbf{Advantages}
    & \textbf{Limitations} & \textbf{Examples} \\
\midrule
Process-based
    & Use mathematical equations to simulate soil formation and
      physical--chemical processes.
    & Provide high accuracy when well-calibrated with specific biological
      data.
    & Can be overcomplicated due to multiple variables.
    & Integrated Critical Zone (ICZ) model \cite{Giannakis2017} \\[6pt]
Empirical
    & Simple models of soil structure useful for data-scarce areas. Rely on
      links between soil and ecosystem services.
    & Efficient for rapid calculation of soil loss (rainfall erodibility,
      slope length, steepness, LULC).
    & Difficulty with dynamic factors (LUC, soil moisture); needs average
      data; overpredictions.
    & Universal Soil Loss Equation (USLE) \cite{Benavidez2018} \\[6pt]
Composite
    & Integrate different approaches: process-based $+$ empirical
      frameworks.
    & Provide a holistic view of soil functions and their contribution to
      ecosystem services.
    & Need accurate validation for joint probability distributions.
    & Life Cycle Assessment (LCA) \cite{Sadhukhan2022} \\[6pt]
Statistical
    & Analyse the complex links in different SES through capturing
      non-linear relationships.
    & Flexibility for high-dimensional data; accurate detection of tail
      dependencies.
    & Low accuracy; need for large datasets; complexity of structure.
    & Vine copulas \cite{Lu2020} \\
\bottomrule
\end{tabularx}
\end{table}

The empirical models yield precise quantification and contribute to the
understanding of soil pollution, although they are challenged by the
ecological soil monitoring limitations associated with the complexity and
variability of different soil types, necessitating further validation. A
prevalent model type for environmental risk assessment estimates ecological
risk factors through contamination indices
\cite{Stojic2019,Mohammad2025,Sun2025,Edo2024,Chon2011,Ferreira2022}. These
methods involve comparing pollutant concentrations in soil samples with
established background levels, where the physical-chemical properties and
contamination levels are examined for ecotoxicity via bioassay techniques. The
summary highlights significant approaches in quantifying ecological indices,
which utilize diverse mathematical frameworks, as outlined in
Table~\ref{tab2}.

\begin{table}[htbp]
\caption{Summary of the major ecological risk indices developed as tools to
    assess harm to soil ecosystems from hazardous
    substances.}\label{tab2}
\begin{tabularx}{\textwidth}{@{} c p{4cm} X p{6cm} @{}}
\toprule
\textbf{No} & \textbf{Name} & \textbf{References} & \textbf{Approach} \\
\midrule
1 & ERI -- Ecological Risk Index
    & \cite{Wang2025,Egbueri2020}
    & Uses contamination factor ($C_f$). \\[4pt]
2 & TRIAD modelling approach (Chemistry, Ecotoxicology, Ecology)
    & \cite{Kim2024}
    & Ecological risk assessment that integrates three lines of evidence:
      chemical, ecotoxicological, and ecological. \\[4pt]
3 & PERI -- Potential Ecological Risk Index
    & \cite{AnsoarRodriguez2025}
    & Uses concentration of HM pollutants and toxicity (harmfulness and
      evaluation of risk at high concentrations). \\[4pt]
4 & $I_{\mathrm{geo}}$ -- Geo-accumulation index
    & \cite{Li2014}
    & Measures HM contamination in soil to assess the increase of HM after
      and before industrial activities by comparing current concentrations to
      background levels. \\[4pt]
5 & MRI -- Modified potential ecological risk index by $I_{\mathrm{geo}}$
    & \cite{Liu2021}
    & Accurately assesses ecological risks of agricultural soils. \\[4pt]
6 & $E_{\mathrm{rf}}$ -- Ecological risk factor
    & \cite{Li2025a}
    & Quantifies by combining contaminant concentrations with their toxicity
      and background levels. \\[4pt]
7 & PLI -- Pollution Load Index
    & \cite{Li2022b,Chukwuonye2025}
    & Determines the degree of pollution in environmental matrices by
      comparing pollutant concentrations in samples to background levels. \\
\bottomrule
\end{tabularx}
\end{table}

Collecting soil data is a fundamental step in soil analysis, requiring a
systematic method to acquire representative samples of soil aggregates. The
process begins with fieldwork, where the data is gathered to facilitate
laboratory analysis. Key characteristics of data sampling include defining the
area and sampling pattern, utilizing advanced equipment to collect multiple
samples from diverse representative categories at target depths, and
thoroughly mixing subsamples to create composite samples. Furthermore, the
metadata associated with soil collection is meticulously recorded, including
detailed field information about the meteorological and topographic context of
the study areas. This comprehensive data collection is essential for
subsequent soil modeling, allowing for a precise evaluation of the
characteristics of soil aggregates, such as texture, structure, biochemical
properties, mineralogical content, and porosity.

% =========================================================================
\section{A Multi-Stage Framework for Soil Pollution Monitoring and SES Modelling}\label{sec4}
% =========================================================================

Challenges and prospects in soil monitoring strategies are designed to ensure the
proper functioning and protection of the soil environment. This complex
approach is visually summarized in Figure~\ref{fig07}, which illustrates the
key components and processes involved in soil monitoring to support
environmental health.

\begin{figure}[htbp]
    \centering
    \includegraphics[width=\textwidth]{Fig_07}
    \caption{Conceptual scheme of SES. Image source: author.}
    \label{fig07}
\end{figure}

The systematic assessment of soil pollution and its impact on ecosystem services necessitates an integrated workflow beginning with the identification of contamination sources through the mapping of hotspot distributions and the laboratory characterization of physicochemical soil properties, including heavy metal concentrations, pH, organic matter, and texture, Table \ref{tab:ses_workflow_v2}. 

\begin{table}[htbp]
\centering
\caption{Workflow Summary for Soil Ecosystem Services (SES) Modelling and Risk Assessment with Supporting Refereces}
\label{tab:ses_workflow_v2}
\begin{tabularx}{\textwidth}{|l|X|p{6.5cm}|X|}
\hline
\textbf{Step} & \textbf{Objective} & \textbf{Methods and Key Components} & \textbf{Supporting Refereces} \\ \hline
1 & \textbf{Detection of pollution sources} & Measurement of initial conditions, hotspot mapping via QGIS, lab analysis of soil properties (HMs, pH, OM, texture), and PCA for spatial distribution assessment. & \cite{Adewumi2024, ElSorogy2025, Feng2024, Lemenkova2021, Mukhopadhyay2020, Wan2019} \\ \hline
2 & \textbf{Modelling pollutant behavior} & Application of empirical, geochemical, and dynamic mass balance models to predict concentration trends and simulate accumulation/attenuation rates. & \cite{Giannakis2017, Koukianaki2025, LindhLemenkova2022b, Mohammad2025, PerezLucas2025, Xu2026} \\ \hline
3 & \textbf{Predictive scenario testing} & Simulation of default, optimistic, and pessimistic scenarios to forecast ecosystem changes and modeling pollutant pathways into the food chain. & \cite{Appel2022, GonzalezRodriguez2011, Hossini2025, Lemenkova2024a, Wang2025, Yasir2025} \\ \hline
4 & \textbf{Evaluating ecosystem risk} & Use of indices ($E_f$, $I_{\mathrm{geo}}$, PERI) to quantify pollution levels and assessment of impacts on soil structure, porosity, and water capacity. & \cite{AnsoarRodriguez2025, Chukwuonye2025, Egbueri2020, Ferreira2022, Kim2024, Li2025a} \\ \hline
5 & \textbf{Multi-factor analysis} & Integration of multiple risk elements for agricultural and forest stands; includes Monte Carlo simulations for uncertainty analysis. & \cite{Lemenkova2025c, Li2025a, Lu2020, Olah2025, Proshad2025, Wang2024b} \\ \hline
6 & \textbf{Prediction and management} & Remote sensing and time series analysis to forecast long-term impacts on health and biodiversity; development of remediation strategies and regulatory policies. & \cite{Fan2025, Lemenkova2025a, Lemenkova2025d, Pan2019, Stojic2019, Xu2025} \\ \hline
\end{tabularx}
\end{table}

By applying Principal Component Analysis (PCA), researchers can delineate the spatial distribution of pollutants to elucidate their environmental impacts and subsequently utilize empirical and geochemical models—specifically dynamic mass balance frameworks—to simulate accumulation and attenuation rates within soil systems. This modeling enables predictive scenario testing across default, optimistic, and pessimistic projections to forecast ecological alterations and trace pollutant translocation from soil to the trophic chain. The quantification of these risks is further refined through the application of ecological indices such as $E_f$, $I_{\text{geo}}$, and PERI, which evaluate the degradation of soil structural integrity, porosity, and hydraulic capacity. To address the complexity of contaminant interactions in agricultural and forest ecosystems, a multi-factor analysis is employed, often reinforced by Monte Carlo simulations to account for inherent uncertainties. Ultimately, the synthesis of remote sensing data and time-series analysis facilitates long-term predictions regarding soil health and biodiversity, providing the data-driven foundation required for formulating remediation strategies and robust environmental management policies.

% =========================================================================
\section{Conclusion}\label{sec5}
% =========================================================================

Soil contamination poses a critical ecological and environmental challenge
globally, impacting both natural ecosystems such as forests and agricultural
sectors like crop production. Effective management and remediation of
contaminated soils hinge on understanding the relationship between pollutants
and soil aggregates, which are characterized by their heterogeneous and
complex structure along with variable solid phase compositions and particle
sizes. This study reviews existing detection methodologies for soil pollution
and proposes strategic frameworks for preventing and addressing contamination.

The literature analysis reveals that current research primarily concentrates
on three main objectives: first, identifying pollution sources in contaminated
soils; second, assessing the adsorption capacity and distribution of
pollutants; and third, examining the phyto-uptake of these pollutants by
plants and their resulting ecotoxicity. Despite significant progress, it is
noted that only a handful of studies in the last decade have provided
comprehensive scoping reviews of soil pollution in soil--plant systems.
Furthermore, successfully managing and remediating soil contamination is
crucial for enhancing our understanding of the interactions between pollutants
and solid soil components within ecosystems. This awareness is vital for
restoring soil health and ensuring the sustainability of agricultural
practices in the face of increasing contamination pressures.

In this study, we examined various sources of soil pollution, including
metals, pesticides, microplastics, per- and polyfluoroalkyl substances
(PFAS), and nutrients. The paper elaborates on the identification methods and
approaches for these contaminants, highlighting the distinctions among
different modeling tools utilized for soil ecosystem modeling. The study
proposes solutions aimed at constructing effective models for identifying and
quantifying soil contaminant concentrations, as well as determining the
origins of these pollutants. Additionally, the project interface integrates an
interdisciplinary analysis of the soil environment, supported by a
comprehensive review of over 100 sources from major bibliographic databases
such as Scopus and Web of Science (WoS). This multifaceted approach
contributes significantly to advancing methodologies for mapping and
identifying soil contaminants.

Major modeling trends for evaluating soil contamination by per- and
polyfluoroalkyl substances (PFAS) were comprehensively reviewed, highlighting
both their strengths and limitations across various approaches. The analysis
demonstrated that current soil monitoring methods effectively support the
traditional framework of soil ecosystem services (SES). This support is
realized through the integration of multi-source datasets, which facilitate
robust evaluations via spectral fingerprints of soil components. Additionally,
the review discussed innovative techniques employed in crop agronomy aimed at
enhancing SES, as well as the performance metrics associated with soil
ecosystem services.

% =========================================================================
% BIBLIOGRAPHY
% =========================================================================
% Each entry below corresponds to a citation key used in the text.
% Replace these \bibitem stubs with your actual BibTeX file or
% insert the formatted references directly.

\begin{thebibliography}{99}

\bibitem{Adewumi2024}
Adewumi, A.J. \& O.D. Ogundele, 2024.
Hidden hazards in urban soils: A meta-analysis review of global heavy metal
contamination (2010--2022), sources and its ecological and health
consequences.
\textit{Sustainable Environment}, 10(1): 2293239.
\url{https://doi.org/10.1080/27658511.2023.2293239}

\bibitem{Akdogan2019}
Akdogan, Z. \& B. Guven, 2019.
Microplastics in the environment: A critical review of current understanding
and identification of future research needs.
\textit{Environmental Pollution}, 254: 113011.
\url{https://doi.org/10.1016/j.envpol.2019.113011}

\bibitem{AliZerrouki2021}
Ali Zerrouki, A. \& M. Melila, 2021.
Evaluation of Soil Contamination by Heavy Metals in the Vicinity of Boucaid
Mine, Ouarsenis (N.O. Algeria).
\textit{Soil and Sediment Contamination: An International Journal},
30(8): 924--942.
\url{https://doi.org/10.1080/15320383.2021.1897084}

\bibitem{Anagnostou2025}
Anagnostou, N., K. Tountopoulos, E. Giorgi, S. Triantafyllaki, P. Rigas \&
P. Samaras, 2025.
Low ppt residue monitoring of persistent organic pollutants in groundwater
samples using direct immersion solid phase micro extraction arrow (DI-SPME)
with gas chromatography-mass spectrometry.
\textit{Microchemical Journal}, 218: 115584.
\url{https://doi.org/10.1016/j.microc.2025.115584}

\bibitem{AnsoarRodriguez2025}
Ansoar-Rodr\'{i}guez, Y., L. Bertrand, C.V. Colombo, G.N. Rimondino,
N. Rivetti, M. d. l. A. Bistoni \& M.V. Am\'{e}, 2025.
Microplastic distribution and potential ecological risk index in a south
american sparsely urbanized river basin.
\textit{Journal of Hazardous Materials Advances}, 18: 100685.
\url{https://doi.org/10.1016/j.hazadv.2025.100685}

\bibitem{Appel2022}
Appel, M., M. Forsthuber, R. Ramos, R. Widhalm, S. Granitzer, M. Uhl,
M. Hengstschl\"{a}ger, T. Stamm \& C. Gundacker, 2022.
The transplacental transfer efficiency of per- and polyfluoroalkyl substances
(PFAS): a first meta-analysis.
\textit{Journal of Toxicology and Environmental Health, Part B},
25(1): 23--42.
\url{https://doi.org/10.1080/10937404.2021.2009946}

\bibitem{Arnesdotter2025}
Arnesdotter, E., C.B.A. Stoffels, W. Alker, A.C. Gutleb \& T. Serchi, 2025.
Per- and polyfluoroalkyl substances (PFAS): immunotoxicity at the primary
sites of exposure.
\textit{Critical Reviews in Toxicology}, 55(4): 484--504.
\url{https://doi.org/10.1080/10408444.2025.2501420}

\bibitem{Atoloye2024}
Atoloye, I., A. Erhunmwunse, R. Ekundayo \& A. Olayinka, 2024.
Compost and mineral fertilizer application affects soil microbial activity
and nutrient mineralization in an ultisol.
\textit{Communications in Soil Science and Plant Analysis},
55(17): 2626--2635.
\url{https://doi.org/10.1080/00103624.2024.2372415}

\bibitem{Baghaie2020}
Baghaie, A.H. \& A. Aghilizefreei, 2020.
Iron enriched green manure can increase wheat Fe concentration in Pb-polluted
soil.
\textit{Soil and Sediment Contamination: An International Journal},
29(7): 721--743.
\url{https://doi.org/10.1080/15320383.2020.1771274}

\bibitem{Benavidez2018}
Benavidez, R., B. Jackson, D. Maxwell \& K. Norton, 2018.
A review of the (revised) universal soil loss equation ((R)USLE).
\textit{Hydrology and Earth System Sciences}, 22(11): 6059--6086.
\url{https://doi.org/10.5194/hess-22-6059-2018}

\bibitem{Bustamante2025}
Bustamante, C.M., P. Ruiz, A.B. Rifat, N. Bravo, J.O. Grimalt \& M. Gar\'{i},
2025.
Methodological approach for a simultaneous determination of persistent and
non-persistent organic pollutants in human blood using gas chromatography and
mass spectrometry techniques.
\textit{Journal of Chromatography A}, 1759: 466235.
\url{https://doi.org/10.1016/j.chroma.2025.466235}

\bibitem{Carr2016}
Carr, S.A., J. Liu \& A.G. Tesoro, 2016.
Transport and fate of microplastic particles in wastewater treatment plants.
\textit{Water Research}, 91: 174--182.
\url{https://doi.org/10.1016/j.watres.2016.01.002}

\bibitem{Chen2020}
Chen, H., Y. Wang, X. Sun, Y. Peng \& L. Xiao, 2020.
Mixing effect of polylactic acid microplastic and straw residue on soil
property and ecological function.
\textit{Chemosphere}, 243: 125271.
\url{https://doi.org/10.1016/j.chemosphere.2019.125271}

\bibitem{Chen2025}
Chen, Z., L.J. Carter, S.A. Banwart \& P. Kay, 2025.
Microplastics in soil--plant systems: Current knowledge, research gaps, and
future directions for agricultural sustainability.
\textit{Agronomy}, 15(7): 1519.
\url{https://doi.org/10.3390/agronomy15071519}

\bibitem{Cheng2025}
Cheng, Y., X. Gao, D. Liu, H. Yu \& Y. Li, 2025.
PFAS in animal by-products in Lanzhou, China, and health risk assessment.
\textit{Food Additives \& Contaminants: Part B}: 1--13.
\url{https://doi.org/10.1080/19393210.2025.2521659}

\bibitem{Chon2011}
Chon, H.-T., J.-S. Lee \& J.-U. Lee, 2011.
Heavy metal contamination of soil, its risk assessment and bioremediation.
\textit{Geosystem Engineering}, 14(4): 191--206.
\url{https://doi.org/10.1080/12269328.2011.10541350}

\bibitem{Choudhary2022}
Choudhary, M., P. Chandra, B. Dixit, V. Nehra, U. Choudhary \& S. Choudhary,
2022.
Plant growth-promoting microbes: Role and prospective in amelioration of salt
stress.
\textit{Communications in Soil Science and Plant Analysis},
53(13): 1692--1711.
\url{https://doi.org/10.1080/00103624.2022.2063316}

\bibitem{Chukwuonye2025}
Chukwuonye, G.N., K. Palawat, R.A. Root et al., 2025.
Using the pollution load index to evaluate rooftop harvested rainwater
metal(loid) contamination in environmental justice communities.
\textit{Environmental Research}, 284: 122187.
\url{https://doi.org/10.1016/j.envres.2025.122187}

\bibitem{Deceuster2012}
Deceuster, J. \& O. Kaufmann, 2012.
Improving the delineation of hydrocarbon-impacted soils and water through
induced polarization (IP) tomographies.
\textit{Journal of Contaminant Hydrology}, 136--137: 25--42.
\url{https://doi.org/10.1016/j.jconhyd.2012.05.003}

\bibitem{EEA2009}
European Environment Agency (EEA), 2009.
Overview of contaminants affecting soil and groundwater in Europe.
[Online chart; modified 11 Sept 2024.]

\bibitem{Edo2024}
Edo, G.I., P.O. Samuel, G.O. Oloni et al., 2024.
Environmental persistence, bioaccumulation, and ecotoxicology of heavy
metals.
\textit{Chemistry and Ecology}, 40(3): 322--349.
\url{https://doi.org/10.1080/02757540.2024.2306839}

\bibitem{Egbueri2020}
Egbueri, J.C., 2020.
Groundwater quality assessment using pollution index of groundwater (PIG),
ecological risk index (ERI) and hierarchical cluster analysis (HCA).
\textit{Groundwater for Sustainable Development}, 10: 100292.
\url{https://doi.org/10.1016/j.gsd.2019.100292}

\bibitem{ElSorogy2025}
El-Sorogy, A.S., K. Al-Kahtany, T. Alharbi, R. Al Hawas \& N. Rikan, 2025.
Geographic Information System and Multivariate Analysis Approach for Mapping
Soil Contamination and Environmental Risk Assessment in Arid Regions.
\textit{Land}, 14(2): 221.
\url{https://doi.org/10.3390/land14020221}

\bibitem{Fan2025}
Fan, X., P. Gao, J. Wang et al., 2025.
Assessment and prediction of soil heavy metal pollution in cotton fields by
multi-source data feature fusion combined with machine learning.
\textit{Industrial Crops and Products}, 233: 121447.
\url{https://doi.org/10.1016/j.indcrop.2025.121447}

\bibitem{Feng2024}
Feng, Y., J. Wang \& Y. Tang, 2024.
Estimation and inversion of soil heavy metal arsenic (As) based on UAV
hyperspectral platform.
\textit{Microchemical Journal}, 207: 112027.
\url{https://doi.org/10.1016/j.microc.2024.112027}

\bibitem{Ferreira2022}
Ferreira, S.L.C., J.B. da Silva, I.F. dos Santos et al., 2022.
Use of pollution indices and ecological risk in the assessment of
contamination from chemical elements in soils and sediments.
\textit{Trends in Environmental Analytical Chemistry}, 35: e00169.
\url{https://doi.org/10.1016/j.teac.2022.e00169}

\bibitem{Gang2004}
Gang, G.Z., L.H. Xia, W.Y. Rong, W.S. Min \& C.G. Dong, 2004.
Suitability of lucerne cultivars to semi-arid conditions in west China.
\textit{New Zealand Journal of Agricultural Research}, 47(1): 51--59.
\url{https://doi.org/10.1080/00288233.2004.9513570}

\bibitem{Giannakis2017}
Giannakis, G.V., N.P. Nikolaidis, J. Valstar et al., 2017.
Chapter Ten -- Integrated Critical Zone Model (1D-ICZ): A Tool for Dynamic
Simulation of Soil Functions and Soil Structure.
\textit{Advances in Agronomy}, 142: 277--314.
\url{https://doi.org/10.1016/bs.agron.2016.10.009}

\bibitem{GonzalezRodriguez2011}
Gonz\'{a}lez-Rodr\'{i}guez, R.M., R. Rial-Otero, B. Cancho-Grande,
C. Gonzalez-Barreiro \& J. Simal-G\'{a}ndara, 2011.
A Review on the Fate of Pesticides during the Processes within the
Food-Production Chain.
\textit{Critical Reviews in Food Science and Nutrition}, 51(2): 99--114.
\url{https://doi.org/10.1080/10408390903432625}

\bibitem{Guan2024}
Guan, R., L. Wu, Y. Li et al., 2024.
Evaluating the Impacts of Fertilization and Rainfall on Multi-Form Phosphorus
Losses from Agricultural Fields.
\textit{Agronomy}, 14(9): 1922.
\url{https://doi.org/10.3390/agronomy14091922}

\bibitem{Guo2021}
Guo, C., F. Chen, J. Li et al., 2021.
Navel orange fine root nutrient content and rhizosphere effects varied with
tree ages and soil depths in a hilly red soil region of China.
\textit{Acta Agriculturae Scandinavica, Section~B}, 71(8): 696--705.
\url{https://doi.org/10.1080/09064710.2021.1940269}

\bibitem{Guo2025}
Guo, P., X. Gu, Z. Li et al., 2025.
A novel almond shell biochar modified with FeS and chitosan as adsorbents for
mitigation of heavy metals from water and soil.
\textit{Separation and Purification Technology}, 360: 130943.
\url{https://doi.org/10.1016/j.seppur.2024.130943}

\bibitem{Hepperly2009}
Hepperly, P., D. Lotter, C.Z. Ulsh, R. Seidel \& C. Reider, 2009.
Compost, manure and synthetic fertilizer influences crop yields, soil
properties, nitrate leaching and crop nutrient content.
\textit{Compost Science \& Utilization}, 17(2): 117--126.
\url{https://doi.org/10.1080/1065657X.2009.10702410}

\bibitem{Hossain2024}
Hossain, M.I., M.A. Khaleque, M.R. Ali et al., 2024.
Development of electrochemical sensors for quick detection of environmental
(soil and water) NPK ions.
\textit{RSC Advances}, 14(13): 9137--9158.
\url{https://doi.org/10.1039/d4ra00034j}

\bibitem{Hossini2025}
Hossini, H., T. Massahi, K. Parnoon \& M. Nouri, 2025.
Per-and polyfluoroalkyl substances (PFAS) in milk and dairy products: a
literature review of the occurrence, contamination sources, and health risks.
\textit{Food Additives \& Contaminants: Part A}, 42(9): 1284--1296.
\url{https://doi.org/10.1080/19440049.2025.2538224}

\bibitem{Hu2020}
Hu, H., J. Zhao, L. Wang et al., 2020.
Synchrotron-based techniques for studying the environmental health effects of
heavy metals: Current status and future perspectives.
\textit{TrAC Trends in Analytical Chemistry}, 122: 115721.
\url{https://doi.org/10.1016/j.trac.2019.115721}

\bibitem{Jing2023}
Jing, X., L. Su, Y. Wang, M. Yu \& X. Xing, 2023.
How do microplastics affect physical properties of silt loam soil under
wetting--drying cycles?
\textit{Agronomy}, 13(3): 844.
\url{https://doi.org/10.3390/agronomy13030844}

\bibitem{Kelsey1968}
Kelsey, J.M. \& E.Z. Arlidge, 1968.
Effects of isobenzan on soil fauna and soil structure.
\textit{New Zealand Journal of Agricultural Research}, 11(2): 245--260.
\url{https://doi.org/10.1080/00288233.1968.10431424}

\bibitem{Kim2024}
Kim, D., J.I. Kwak, T.-Y. Lee et al., 2024.
Triad method to assess ecological risks of contaminated soils in abandoned
mine sites.
\textit{Journal of Hazardous Materials}, 461: 132535.
\url{https://doi.org/10.1016/j.jhazmat.2023.132535}

\bibitem{Klaucol2013}
Klau\v{c}o, M., B. Gregorov\'{a}, U. Stankov, V. Markovi\'{c} \&
P. Lemenkova, 2013.
Determination of ecological significance based on geostatistical assessment:
a case study from the Slovak Natura 2000 protected area.
\textit{Open Geosciences}, 5(1): 28--42.
\url{https://doi.org/10.2478/s13533-012-0120-0}

\bibitem{Klaucol2017}
Klau\v{c}o, M., B. Gregorov\'{a}, P. Koleda, U. Stankov, V. Markovic \&
P. Lemenkova, 2017.
Land Planning as a Support for Sustainable Development Based on Tourism: A
Case Study of Slovak Rural Region.
\textit{Environmental Engineering and Management Journal},
16(2): 449--458.
\url{https://doi.org/10.30638/eemj.2017.045}

\bibitem{Koukianaki2025}
Koukianaki, E.A., M.A. Lilli, D. Efstathiou et al., 2025.
Modeling soil functions of forested ecosystems.
\textit{Journal of Environmental Management}, 386: 125636.
\url{https://doi.org/10.1016/j.jenvman.2025.125636}

\bibitem{Lemenkova2020}
Lemenkova, P., 2020.
GRASS GIS for topographic and geophysical mapping of the Peru-Chile Trench.
\textit{Forum geografic}, 19(2): 143--157.
\url{https://doi.org/10.5775/fg.2020.009.d}

\bibitem{Lemenkova2021}
Lemenkova, P., 2021.
Geophysical Mapping of Ghana Using Advanced Cartographic Tool GMT\@.
\textit{Kartografija i Geoinformacije}, 20(36): 16--37.
\url{https://doi.org/10.32909/kg.20.36.2}

\bibitem{Lemenkova2022}
Lemenkova, P., 2022.
GRASS GIS Scripts for Satellite Image Analysis by Raster Calculations Using
Modules r.mapcalc, d.rgb, r.slope.aspect.
\textit{Tehnicki Vjesnik}, 29(6): 1956--1963.
\url{https://doi.org/10.17559/TV-20220322091846}

\bibitem{Lemenkova2023}
Lemenkova, P., 2023.
Using open-source software GRASS GIS for analysis of the environmental
patterns in Lake Chad, Central Africa.
\textit{Die Bodenkultur: Journal of Land Management, Food and Environment},
74(1): 49--64.
\url{https://doi.org/10.2478/boku-2023-0005}

\bibitem{Lemenkova2024a}
Lemenkova, P., 2024a.
Artificial Intelligence for Computational Remote Sensing: Quantifying
Patterns of Land Cover Types around Cheetham Wetlands, Port Phillip Bay,
Australia.
\textit{Journal of Marine Science and Engineering}, 12(8): 1279.
\url{https://doi.org/10.3390/jmse12081279}

\bibitem{Lemenkova2024b}
Lemenkova, P., 2024b.
[See author's publication list for full reference.]

\bibitem{Lemenkova2025a}
Lemenkova, P., 2025a.
Machine Learning Algorithms of Remote Sensing Data Processing for Mapping
Changes in Land Cover Types over Central Apennines, Italy.
\textit{Journal of Imaging}, 11(5): 153.
\url{https://doi.org/10.3390/jimaging11050153}

\bibitem{Lemenkova2025b}
Lemenkova, P., 2025b.
Reclassification Scheme for Image Analysis in GRASS GIS Using Gradient
Boosting Algorithm: A Case of Djibouti, East Africa.
\textit{Journal of Imaging}, 11(8): 249.
\url{https://doi.org/10.3390/jimaging11080249}

\bibitem{Lemenkova2025c}
Lemenkova, P., 2025c.
On forest hydrology in South Tyrol: Processes, modelling and assessment.
\textit{Agriculture and Forestry}, 71(2): 7--22.
\url{https://doi.org/10.17707/AgricultForest.71.2.01}

\bibitem{Lemenkova2025d}
Lemenkova, P., 2025d.
Improving Bimonthly Landscape Monitoring in Morocco, North Africa, by
Integrating Machine Learning with GRASS GIS\@.
\textit{Geomatics}, 5(1): 5.
\url{https://doi.org/10.3390/geomatics5010005}

\bibitem{Li2014}
Li, Z., Z. Ma, T.J. van der Kuijp, Z. Yuan \& L. Huang, 2014.
A review of soil heavy metal pollution from mines in China: Pollution and
health risk assessment.
\textit{Science of The Total Environment}, 468--469: 843--853.
\url{https://doi.org/10.1016/j.scitotenv.2013.08.090}

\bibitem{Li2019}
Li, C., K. Zhou, W. Qin et al., 2019.
A review on heavy metals contamination in soil: Effects, sources, and
remediation techniques.
\textit{Soil and Sediment Contamination: An International Journal},
28(4): 380--394.
\url{https://doi.org/10.1080/15320383.2019.1592108}

\bibitem{Li2022b}
Li, F.-j., H.-w. Yang, R. Ayyamperumal \& Y. Liu, 2022.
Pollution, sources, and human health risk assessment of heavy metals in urban
areas around industrialization and urbanization---northwest China.
\textit{Chemosphere}, 308: 136396.
\url{https://doi.org/10.1016/j.chemosphere.2022.136396}

\bibitem{Li2025a}
Li, X., D. Ding, X. Wang et al., 2025.
Integration of self-organizing map and Monte Carlo simulation for ecological
risk prediction of heavy metal attenuation in groundwater.
\textit{Ecotoxicology and Environmental Safety}, 302: 118761.
\url{https://doi.org/10.1016/j.ecoenv.2025.118761}

\bibitem{LindhLemenkova2022a}
Lindh, P. \& P. Lemenkova, 2022a.
Soil contamination from heavy metals and persistent organic pollutants
(PAH, PCB and HCB) in the coastal area of V\"{a}sternorrland, Sweden.
\textit{Gospodarka Surowcami Mineralnymi -- Mineral Resources Management},
38(2): 147--168.
\url{https://doi.org/10.24425/gsm.2022.141662}

\bibitem{LindhLemenkova2022b}
Lindh, P. \& P. Lemenkova, 2022b.
Geochemical tests to study the effects of cement ratio on potassium and TBT
leaching and the pH of the marine sediments from the Kattegat Strait, Port of
Gothenburg, Sweden.
\textit{Baltica}, 35(1): 47--59.
\url{https://doi.org/10.5200/baltica.2022.1.4}

\bibitem{LindhLemenkova2022d}
Lindh, P. \& P. Lemenkova, 2022d.
Simplex Lattice Design and X-ray Diffraction for Analysis of Soil Structure.
\textit{Electronics}, 11(22): 3726.
\url{https://doi.org/10.3390/electronics11223726}

\bibitem{LindhLemenkova2023a}
Lindh, P. \& P. Lemenkova, 2023a.
Effects of Water--Binder Ratio on Strength and Seismic Behavior of Stabilized
Soil from Kongshavn, Port of Oslo.
\textit{Sustainability}, 15(15): 12016.
\url{https://doi.org/10.3390/su151512016}

\bibitem{LindhLemenkova2023b}
Lindh, P. \& P. Lemenkova, 2023b.
Optimized Workflow Framework in Construction Projects to Control the
Environmental Properties of Soil.
\textit{Algorithms}, 16(6): 303.
\url{https://doi.org/10.3390/a16060303}

\bibitem{LindhLemenkova2023c}
Lindh, P. \& P. Lemenkova, 2023c.
Hardening Accelerators as Strength-Enhancing Admixture Solutions for Soil
Stabilization.
\textit{Slovak Journal of Civil Engineering}, 31(1): 10--21.
\url{https://doi.org/10.2478/sjce-2023-0002}

\bibitem{LindhLemenkova2023d}
Lindh, P. \& P. Lemenkova, 2023d.
Seismic monitoring of strength in stabilized foundations.
\textit{Journal of the Mechanical Behavior of Materials}, 32(1): 20220290.
\url{https://doi.org/10.1515/jmbm-2022-0290}

\bibitem{LindhLemenkova2023e}
Lindh, P. \& P. Lemenkova, 2023e.
Behaviour of Moraine Soils Stabilised with OPC, GGBFS and Hydrated Lime.
\textit{Archives of Mining Sciences}, 68(2): 319--334.
\url{https://doi.org/10.24425/ams.2023.146182}

\bibitem{Liu2021}
Liu, X., S. Chen, X. Yan et al., 2021.
Evaluation of potential ecological risks in potential toxic elements
contaminated agricultural soils.
\textit{Journal of Environmental Management}, 300: 113679.
\url{https://doi.org/10.1016/j.jenvman.2021.113679}

\bibitem{Lu2020}
L\"{u}, T.-J., X.-S. Tang, D.-Q. Li \& X.-H. Qi, 2020.
Modeling multivariate distribution of multiple soil parameters using vine
copula model.
\textit{Computers and Geotechnics}, 118: 103340.
\url{https://doi.org/10.1016/j.compgeo.2019.103340}

\bibitem{Lu2023}
Lu, X., F. Li, W. Yang, P. Zhu \& S. Lv, 2023.
Quantitative analysis of heavy metals in soil by X-ray fluorescence with
improved variable selection strategy and Bayesian optimized support vector
regression.
\textit{Chemometrics and Intelligent Laboratory Systems}, 238: 104842.
\url{https://doi.org/10.1016/j.chemolab.2023.104842}

\bibitem{Luo2017}
Luo, Y., M.B. Benke \& X. Hao, 2017.
Nutrient uptake and leaching from soil amended with cattle manure and
nitrapyrin.
\textit{Communications in Soil Science and Plant Analysis},
48(12): 1438--1454.
\url{https://doi.org/10.1080/00103624.2017.1373786}

\bibitem{Ma2024}
Ma, D., W. Teng, Y.-T. Mo et al., 2024.
Effects of nitrogen, phosphorus, and potassium fertilization on plant growth
and element levels in \textit{Erythropalum scandens}.
\textit{Journal of Plant Nutrition}, 47(1): 82--96.
\url{https://doi.org/10.1080/01904167.2023.2262504}

\bibitem{Marcelino2024}
Marcelino, S.M., P.D. Gaspar, A. do Pa\c{c}o et al., 2024.
Agricultural Practices for Biodiversity Enhancement: Evidence and
Recommendations for the Viticultural Sector.
\textit{AgriEngineering}, 6(2): 1175--1194.
\url{https://doi.org/10.3390/agriengineering6020067}

\bibitem{Mohammad2025}
Mohammad, S.J., Y.E. Ling, K.A. Halim, B.S. Sani \& N.I. Abdullahi, 2025.
Heavy metal pollution and transformation in soil: a comprehensive review of
natural bioremediation strategies.
\textit{Journal of Umm Al-Qura University for Applied Sciences},
11(3): 528--544.
\url{https://doi.org/10.1007/s43994-025-00241-6}

\bibitem{Mozaffari2026}
Mozaffari, H., A.A. Moosavi, J.M. Demat\^{e} \& W. Cornelis, 2026.
A novel and simple method for accurate prediction of soil particle-size
distribution from limited soil texture data.
\textit{Soil and Tillage Research}, 256: 106858.
\url{https://doi.org/10.1016/j.still.2025.106858}

\bibitem{Mukhopadhyay2020}
Mukhopadhyay, S., S. Chakraborty, P.B.S. Bhadoria, B. Li \&
D.C. Weindorf, 2020.
Assessment of heavy metal and soil organic carbon by portable X-ray
fluorescence spectrometry and NixPro\texttrademark\ sensor in landfill soils
of India.
\textit{Geoderma Regional}, 20: e00249.
\url{https://doi.org/10.1016/j.geodrs.2019.e00249}

\bibitem{Nauanova2025}
Nauanova, A., A. Kiyas, S. Kenzhegulova et al., 2025.
The influence of crop rotations, fertilizer and pesticide application on soil
microbial diversity, community composition, and wheat yield.
\textit{Cogent Food \& Agriculture}, 11(1): 2529367.
\url{https://doi.org/10.1080/23311932.2025.2529367}

\bibitem{Olah2025}
Ol\'{a}h, P. \& P. G\"{o}r\"{o}g, 2025.
Integrating Soil Parameter Uncertainty into Slope Stability Analysis: A Case
Study of an Open Pit Mine in Hungary.
\textit{Geosciences}, 15(6): 222.
\url{https://doi.org/10.3390/geosciences15060222}

\bibitem{Pal2024}
Pal, A., S.K. Dubey, S. Goel \& P.K. Kalita, 2024.
Portable sensors in precision agriculture: assessing advances and challenges
in soil nutrient determination.
\textit{TrAC Trends in Analytical Chemistry}, 180: 117981.
\url{https://doi.org/10.1016/j.trac.2024.117981}

\bibitem{Pan2019}
Pan, G., X. Miao, L. Bi et al., 2019.
Modified Local Soil (MLS) Technology for Harmful Algal Bloom Control,
Sediment Remediation, and Ecological Restoration.
\textit{Water}, 11(6): 1123.
\url{https://doi.org/10.3390/w11061123}

\bibitem{PerezLucas2025}
P\'{e}rez-Lucas, G., A. Mart\'{i}nez-Zapata \& S. Navarro, 2025.
Valorization of agro-industrial wastes as organic amendments to reduce
herbicide leaching into soil.
\textit{Journal of Xenobiotics}, 15(4): 100.
\url{https://doi.org/10.3390/jox15040100}

\bibitem{Polakowski2023}
Polakowski, C., A. Mak\'{o}, A. Sochan et al., 2023.
Recommendations for soil sample preparation, pretreatment, and data
conversion for texture classification in laser diffraction particle size
analysis.
\textit{Geoderma}, 430: 116358.
\url{https://doi.org/10.1016/j.geoderma.2023.116358}

\bibitem{Ponder2002}
Ponder, F., 2002.
Implications of silvicultural pesticides on forest soil animals, insects,
fungi, and other organisms and ramifications on soil fertility and health.
\textit{Communications in Soil Science and Plant Analysis},
33(11--12): 1927--1940.
\url{https://doi.org/10.1081/CSS-120004833}

\bibitem{Proshad2025}
Proshad, R., S.M. Asharaful Abedin Asha, R. Tan et al., 2025.
Machine learning models with innovative outlier detection techniques for
predicting heavy metal contamination in soils.
\textit{Journal of Hazardous Materials}, 481: 136536.
\url{https://doi.org/10.1016/j.jhazmat.2024.136536}

\bibitem{Racic2025}
Ra\v{c}i\'{c}, N., S. Ru\v{z}i\v{c}i\'{c}, G. Pehnec et al., 2025.
Linking Atmospheric and Soil Contamination: A Comparative Study of PAHs and
Metals in PM10 and Surface Soil near Urban Monitoring Stations.
\textit{Toxics}, 13(10): 866.
\url{https://doi.org/10.3390/toxics13100866}

\bibitem{Rahman2014}
Rahman, A., S. Rahman \& L. Cihacek, 2014.
Influence of soil pH in vegetative filter strips for reducing soluble nutrient
transport.
\textit{Environmental Technology}, 35(14): 1744--1752.
\url{https://doi.org/10.1080/09593330.2014.881421}

\bibitem{Rasafi2021}
Rasafi, T.E., M. Nouri \& A. Haddioui, 2021.
Metals in mine wastes: environmental pollution and soil remediation
approaches---a review.
\textit{Geosystem Engineering}, 24(3): 157--172.
\url{https://doi.org/10.1080/12269328.2017.1400474}

\bibitem{Rosskopfova2012}
Rosskopfov\'{a}, O., M. Galambos, J. Ometakov\'{a}, M. \v{C}aplo\v{c}i\v{c}ov\'{a}
\& P. Rajec, 2012.
Study of sorption processes of copper on synthetic hydroxyapatite.
\textit{Journal of Radioanalytical and Nuclear Chemistry}, 293: 641--647.
\url{https://doi.org/10.1007/s10967-012-1711-4}

\bibitem{Sadhukhan2022}
Sadhukhan, J., 2022.
Net zero electricity systems in global economies by life cycle assessment
(LCA).
\textit{Renewable Energy}, 184: 960--974.
\url{https://doi.org/10.1016/j.renene.2021.12.024}

\bibitem{Sandu2025}
Sandu, M.A., M. Preda, V. Tanase et al., 2025.
Trends in Polychlorinated Biphenyl Contamination in Bucharest's Urban Soils:
A Two-Decade Perspective (2002--2022).
\textit{Processes}, 13(5): 1357.
\url{https://doi.org/10.3390/pr13051357}

\bibitem{Shen2024}
Shen, Y., L. Wang, Y. Ding et al., 2024.
Trends in the analysis and exploration of per- and polyfluoroalkyl substances
(PFAS) in environmental matrices: A review.
\textit{Critical Reviews in Analytical Chemistry}, 54(8): 3171--3195.
\url{https://doi.org/10.1080/10408347.2023.2231535}

\bibitem{Stojic2019}
Stoji\'{c}, N., S. \v{S}trbac \& D. Proki\'{c}, 2019.
Soil Pollution and Remediation.
In: Hussain, C. (ed.) \textit{Handbook of Environmental Materials
Management}. Springer, Cham.
\url{https://doi.org/10.1007/978-3-319-73645-7_81}

\bibitem{Sun2025}
Sun, X., L. Zhang, Y. Gu et al., 2025.
Nutrient-element-mediated alleviation of cadmium stress in plants:
Mechanistic insights and practical implications.
\textit{Plants}, 14(19): 3081.
\url{https://doi.org/10.3390/plants14193081}

\bibitem{Sutherland2010}
Sutherland, R.A., 2010.
BCR\textsuperscript{\textregistered}-701: A review of 10-years of sequential
extraction analyses.
\textit{Analytica Chimica Acta}, 680(1): 10--20.
\url{https://doi.org/10.1016/j.aca.2010.09.016}

\bibitem{TeranGomez2025}
Ter\'{a}n-G\'{o}mez, V.F., A.M. Buitrago-Ram\'{i}rez, A.F. Echeverri-S\'{a}nchez,
A. Figueroa-Casas \& J.A. Benavides-Bola\~{n}os, 2025.
Integrating AHP and GIS for Sustainable Surface Water Planning.
\textit{Sustainability}, 17(9): 4130.
\url{https://doi.org/10.3390/su17094130}

\bibitem{Toumi2025}
Toumi, K., A. Arbi, N. Soudani et al., 2025.
Ecotoxicological risk assessment and monitoring of pesticide residues in soil,
surface water, and groundwater in northwestern Tunisia.
\textit{Water}, 17(16): 2387.
\url{https://doi.org/10.3390/w17162387}

\bibitem{Vishwanath2025}
Vishwanath, Kumar, S., T.J. Purakayastha et al., 2025.
Five decades of nutrient management enhances soil microbial activity,
resistance and resilience against drought stress.
\textit{Communications in Soil Science and Plant Analysis},
56(14): 2189--2206.
\url{https://doi.org/10.1080/00103624.2025.2494849}

\bibitem{Wan2019}
Wan, M., W. Hu, M. Qu et al., 2019.
Application of arc emission spectrometry and portable X-ray fluorescence
spectrometry to rapid risk assessment of heavy metals in agricultural soils.
\textit{Ecological Indicators}, 101: 583--594.
\url{https://doi.org/10.1016/j.ecolind.2019.01.069}

\bibitem{Wang2020}
Wang, Z., Y. Luo, C. Zheng, C. An \& Z. Mi, 2020.
Spatial distribution, source identification, and risk assessment of heavy
metals in the soils from a mining region.
\textit{Human and Ecological Risk Assessment: An International Journal},
27(5): 1276--1295.
\url{https://doi.org/10.1080/10807039.2020.1821350}

\bibitem{Wang2022}
Wang, F., Q. Wang, C.A. Adams, Y. Sun \& S. Zhang, 2022.
Effects of microplastics on soil properties: Current knowledge and future
perspectives.
\textit{Journal of Hazardous Materials}, 424: 127531.
\url{https://doi.org/10.1016/j.jhazmat.2021.127531}

\bibitem{Wang2024a}
Wang, M., X. Jiang, Z. Wei, L. Wang, J. Song \& P. Cen, 2024.
Enhanced cadmium adsorption dynamics in water and soil by polystyrene
microplastics and biochar.
\textit{Nanomaterials}, 14(13): 1067.
\url{https://doi.org/10.3390/nano14131067}

\bibitem{Wang2024b}
Wang, B., S. Xin, D. Jin, L. Zhang, J. Wu \& H. Guo, 2024.
A New Porosity Evaluation Method Based on a Statistical Methodology for
Granular Material.
\textit{Applied Sciences}, 14(16): 7379.
\url{https://doi.org/10.3390/app14167379}

\bibitem{Wang2025}
Wang, F., Y. Wu, Y. Zhang et al., 2025.
Research on ecosystem service value and landscape ecological risk prediction
and zoning.
\textit{Ecological Modelling}, 507: 111173.
\url{https://doi.org/10.1016/j.ecolmodel.2025.111173}

\bibitem{Xu2023}
Xu, B., G. Yang, A. Lehmann, S. Riedel \& M.C. Rillig, 2023.
Effects of perfluoroalkyl and polyfluoroalkyl substances (PFAS) on soil
structure and function.
\textit{Soil Ecology Letters}, 5(1): 108--117.
\url{https://doi.org/10.1007/s42832-022-0143-5}

\bibitem{Xu2025}
Xu, S., J. Lin, H. Luo et al., 2025.
Technologies for the remediation of nitrogen pollution and advances in the
application of metal--phenolic networks.
\textit{Processes}, 13(7): 2131.
\url{https://doi.org/10.3390/pr13072131}

\bibitem{Xu2026}
Xu, Y., S. Xu, Y. Cheng et al., 2026.
Distinct molecular complexing mechanisms of smoke dissolved organic matters
toward heavy metals: New insights based on FT-ICR-MS, two dimensional
fluorescence EEM and FTIR\@.
\textit{Fuel}, 405: 136646.
\url{https://doi.org/10.1016/j.fuel.2025.136646}

\bibitem{Yasir2025}
Yasir, M., A. Hossain \& A. Pratap-Singh, 2025.
Pesticide degradation: Impacts on soil fertility and nutrient cycling.
\textit{Environments}, 12(8): 272.
\url{https://doi.org/10.3390/environments12080272}

\bibitem{Zireeni2023}
Zireeni, Y., D.L. Jones \& D.R. Chadwick, 2023.
Influence of slurry acidification with H$_2$SO$_4$ on soil pH, N, P, S, and
C dynamics: incubation experiment.
\textit{Environmental Advances}, 14: 100447.
\url{https://doi.org/10.1016/j.envadv.2023.100447}

\bibitem{Zhou2024}
Zhou, J., H. Xu, Y. Xiang \& J. Wu, 2024.
Effects of microplastics pollution on plant and soil phosphorus: A
meta-analysis.
\textit{Journal of Hazardous Materials}, 461: 132705.
\url{https://doi.org/10.1016/j.jhazmat.2023.132705}

\end{thebibliography}

\end{document}
